\documentclass[11pt]{article}
\usepackage[margin=25mm]{geometry}
\usepackage{fontspec,unicode-math}
\setmainfont{Latin Modern Roman}
\setmathfont{Latin Modern Math}
\usepackage{amsmath,amssymb,mathtools}
\usepackage{xurl,hyperref}
\hypersetup{colorlinks=true,urlcolor=blue,linkcolor=blue}
\setcounter{secnumdepth}{0}
\setlength{\parindent}{0pt}
\setlength{\parskip}{5pt}
\setlength{\emergencystretch}{3em}
\newcommand{\sref}[2]{\hyperlink{source:#1:#2}{[S#2]}}
\newcommand{\cataloguescope}[1]{\paragraph{Source question.}#1}
\newcommand{\floor}[1]{\left\lfloor#1\right\rfloor}
\begin{document}
\section{Efficient Moment Relaxations for Polynomial Optimization}\label{statement}

\subsection{Definitions and mathematical statement}

\paragraph{Shared notation.}Standard number systems are denoted by $\mathbb N,\mathbb Z,\mathbb Q,\mathbb R,\mathbb C$; index ranges, computational models and specialized notation follow the cited problem source and the definitions below.


A polynomial is a sum of squares (SOS) if it is $\sum_j q_j^2$ for polynomials $q_j$. Moment-SOS hierarchies impose positivity through increasingly large truncated moment matrices and corresponding degree-bounded SOS certificates.

A semialgebraic set is described by finitely many polynomial equalities and inequalities and finite Boolean combinations of these conditions. A homeomorphism is a continuous bijection with continuous inverse.
\paragraph{Statement recorded in the supplied source.}
Can moment-relaxation methods solve polynomial optimization over closed basic semialgebraic sets efficiently, including computing the infimum and the points where it is attained? \sref{statement}{2}
\paragraph{Scope qualification.}
The oracle/encoding model, approximation accuracy, boundedness, and meaning of efficient are not specified; infimum attainment is not assumed.
\paragraph{Source status.}
The supplied catalogue labels this record open. Its source-status metadata was checked against the supplied snapshot on 2026-10-11. This import does not independently establish current open status or mathematical correctness. The catalogue literature assessment is dated 2026-08-21.
\subsection{Short English statement}
Obtain efficient moment-based global optimization over the stated semialgebraic feasible sets, including minimizers.
\subsection{Sources}
\begin{itemize}
\item[\hypertarget{source:statement:1}{S1}] Ulam AI and the UnsolvedMath Contributors, UnsolvedMath: A Curated Collection of Open Mathematics Problems, catalogue entry. Catalogue curation and metadata: CC BY 4.0. The supplied source question is retained; definitions and the English summary are editorial additions.
\url{https://huggingface.co/datasets/ulamai/UnsolvedMath}.
\item[\hypertarget{source:statement:2}{S2}] Real Algebraic Geometry With A View Toward Systems Control and Free Positivity (2015), Oberwolfach Reports, DOI 10.4171/owr/2014/17, pp. 977--1045
\url{https://doi.org/10.4171/owr/2014/17}.
\item[\hypertarget{source:statement:3}{S3}] Mai--Lasserre--Magron--Wang, arXiv
\url{https://arxiv.org/abs/2012.08873}.
\textit{Underlying Oberwolfach Reports material retains its source-specific licence, including CC BY-SA 4.0 where applicable; see THIRD\_PARTY\_NOTICES.md.}
\end{itemize}
\paragraph{Preparation and review.}
Original question attributed to the cited source and UnsolvedMath contributors. Codex (GPT-6; exact serving revision not exposed) assisted with definitions, English summaries, source inspection and formatting. No new human mathematical authorship or independent proof verification is asserted. Catalogue identity checked against the complete supplied registry and existing repository statements; known flagged repetitions and obvious identical conjectures excluded. Full mathematical-equivalence/literature review has not been performed. Statement preparation and final review followed contribution prompts C and D; automated checks verify structure and identity, not open status.
\label{end:statement}
\end{document}
